\SourceLeaf{388}%
man doch einen grossen Theil des Inhalts nur für Dissertationen erkennen, in denen elementarische Gegenstände mit mehr Breite als Tiefe abgehandelt werden. Es würde jedoch unpassend sein, diese Arbeiten. die einer längst vergangenen Zeit angehören, jetzt noch einer speciellen Kritik zu unterwerfen.

Weit gehaltvoller erscheint dagegen die neue Fortsetzung, worin der Militär, der Geschichtsforscher, der Geograph reichen Stoff zur Belehrung antreffen. Wir nennen hier nur die Darstellung der Schlacht bei Marengo, die Geschichte des Feldzugs in Deutschland im Jahre 1800 (welche beinahe den ganzen fünften Band ausfüllt), die militärische Beschreibung des Flussgebiets der Donau, alles durch eine grosse Menge von Karten und Planen erläutert; die Verhandlungen einer besonders dazu niedergesetzten Commission über die zweckmässigste Art der Terraindarstellung, worin dieser Gegenstand vielseitig erwogen und durch eine beträchtliche Anzahl von Probezeichnungen nach verschiedenen Methoden versinnlicht wird. Nicht ohne Interesse wird man in der Correspondenz des Grafen \textsc{de Gisors} mit seinem Vater dem Herzog \textsc{de Belle-Isle} die Unterredungen lesen, welche ersterer mit \textsc{Friedrich} dem Zweiten ein Jahr vor dem Ausbruche des siebenjährigen Krieges über militärische Gegenstände hatte; imgleichen eine Reihe von bisher ungedruckten zum Theil eigenhändigen, auch mit einem Facsimile begleiteten Briefen \textsc{Ludwig} XIV., wenn gleich nicht alle Leser sie von dem Standpunkt betrachten können, auf welchen die Herausgeber die französischen Leser stellen wollen, indem sie in der Einleitung dazu bemerken: \textsc{Montesquieu} \textit{regarde comme le devoir de tout écrivain homme de bien de contribuer, autant qu’il est en lui, à donner à ses concitoyens des raisons d’aimer ceux à qui ils doivent obéir. Rien n’entre mieux dans cette noble pensée de} \textsc{Montesquieu}, \textit{que la publication de ces lettres et de celles qui pourront les suivre, et tout le monde reconnaîtra dans les successeurs du grand roi tout ce que son coeur avait de paternel et son âme d’héroïque.} Endlich dürfen wir nicht mit Stillschweigen übergehen die Nachrichten, welche im 3. und 4. Bande über die neue grosse Karte von Frankreich gegeben werden, deren Ausführung durch eine königliche Ordonnanz vom 6. August 1817 befohlen wurde. Man wollte Anfangs die Aufnahme in dem Maassstabe von 1 zu 10000 und den Stich in dem Maassstabe von 1 zu 50000 ausführen, wobei die Anzahl aller Blätter auf 611, jedes 800 Millimeter breit und 500 Millimeter hoch, angeschlagen wurde, und glaubte die ganze Arbeit in 20 Jahren vollenden zu können. Man liess jedoch diesen Plan bald fahren, und beschränkte den Maassstab für die Auf%
